\begin{abstract}

We present the instrument design and initial results for the Galactic Radio Explorer (GReX), an all-sky monitor for exceptionally bright transients in the radio sky.
This instrument builds on the success of STARE2 to search for fast radio bursts (FRBs) from the Milky Way and its satellites.
This instrument has deployments across the globe, with wide sky coverage and searching down to $32\,\mu\text{s}$ time resolution, enabling the discovery of new super giant pulses.
Presented here are the details of the hardware and software design of the instrument, performance in sensitivity and other key metrics, and experience in building a global-scale, low-cost experiment.
We follow this discussion with experimental results on validation of the sensitivity via hydrogen-line measurements. We then update the rate of Galactic FRBs based on a non-detection in $0.5\,\text{sky}\,\text{yr}$ of exposure, with twice the sensitivity of STARE2. 
Our results suggest FRB-like events are even rarer than initially implied by the detection of a MJy burst from SGR J1935+2154 in April 2020.

\end{abstract}


\keywords{\uat{Fast radio bursts}{2008} --- \uat{Instrumentation}{799} --- \uat{Magnetars}{992}}

